\section{Technical Supplement: Mathematical Foundations of OPH}\label{technical-supplement-mathematical-foundations-of-observer-patch-holography}

\begin{quote}
OPH is an observer-centric reconstruction program for fundamental physics. In the current implementation it uses two external inputs (pixel area and screen capacity) together with structural axioms and explicit regulator/scaling premises. This supplement collects the mathematical ingredients for the conditional Lorentz/null-modular/Einstein branch, conditional compact gauge reconstruction, and the related bosonic quantitative program.
\end{quote}

\subsection{For Theoretical Physicists}\label{for-theoretical-physicists}

This document collects mathematical derivations and stated premises supporting the OPH program. The Lorentz/null-modular/Einstein and compact-gauge statements are scaling-limit / categorical results under explicit regularity hypotheses; these hypotheses are target properties of the intended EFT phase and candidate UV completions, not fixed-cutoff matrix-algebra theorems.

\textbf{Scope convention.} Statements about Lorentz recovery, modular flow, and the Einstein branch are scaling-limit / EFT statements rather than literal fixed-cutoff type-I claims. Statements about gauge reconstruction refer to the refinement-stable directed colimit of fixed-cutoff sector categories, not to the finite block list at one regulator scale.

\begin{center}\rule{0.5\linewidth}{0.5pt}\end{center}

\section{1. Core Axioms and Definitions}\label{1-core-axioms-and-definitions}

\subsection{1.1 Canonical Five-Axiom Package}\label{11-the-axiom-set}

This supplement uses the same authoritative OPH premise ledger as the compact note and the main manuscript. The canonical axioms are:

\textbf{Screen Net:} Physical reality is encoded on a horizon screen \(S^2\) carrying a net of von Neumann algebras \(P \mapsto \mathcal{A}(P)\) for patches \(P \subset S^2\).

\textbf{Overlap Consistency:} For overlapping patches \(P_1 \cap P_2 \neq \emptyset\), the restrictions of local states must agree: \[\omega_1|_{\mathcal{A}(P_1 \cap P_2)} = \omega_2|_{\mathcal{A}(P_1 \cap P_2)}\]

\textbf{Local MaxEnt and Refinement Stability:} At the UV scale the realized state maximizes entropy subject to a finite family of gauge-invariant local constraints, and the resulting family is stable under coarse-graining.

\textbf{Recoverable Generalized Entropy:} A generalized entropy functional exists: \[S_{\text{gen}}(C) = \frac{A(\partial C)}{4G} + S_{\text{bulk}}(C)\] satisfying quantum focusing and the recoverability structure used in collar and null-strip arguments.

\textbf{Minimal Admissible Realization (MAR):} On the admissible low-energy branch, the realized sector package is the lexicographically minimal one under the complexity vector \(C(\mathfrak S)\). MAR enters the gauge-selection branch; the Lorentz/null-modular and Einstein discussions below mainly use the first four axioms plus the technical premises stated where needed.

\subsection{1.2 Technical Premises and Legacy Regulator Names}\label{12-regulator-premises}

The supplement continues to use historical labels such as R0, R1, and N1--N3 because later derivations cite them directly. In the present version, R0/R1 are not extra microphysical axioms. They are read-along names for the fixed-cutoff realization of the canonical screen-net and overlap-gluing data.

\textbf{R0 (finite type-I regulator presentation):} On a fixed UV cellulation, each cell carries finite local data and can be Hilbertized on a finite-dimensional space \(\tilde{\mathcal H}_i\). For any finite union \(R\) of cells, the extended algebra before quotienting by overlap redundancy is
\[
\widetilde{\mathcal A}(R)=\mathcal B(\tilde{\mathcal H}_R),
\qquad
\tilde{\mathcal H}_R=\bigotimes_{i\subset R}\tilde{\mathcal H}_i.
\]

\textbf{R1 (boundary fixed-point realization of overlap gluing):} The redundancy in overlap identifications acts only on cut data. On each finite-dimensional chart, its unitary image has compact closure, so for a region \(R \subset S^2\) one may represent it by a compact boundary group \(G_{\partial R}\) acting on \(\tilde{\mathcal H}_R\), with physical algebra
\[
\mathcal{A}(R) = \mathcal{B}(\tilde{\mathcal{H}}_R)^{G_{\partial R}}.
\]

\textbf{Read-along convention.} Older references to "A3" mean Recoverable Generalized Entropy; older references to "A4" mean the recoverability clause used together with the collar/null-strip control hypotheses stated explicitly below. The remaining nontrivial burden is the controlled refinement limit and the later geometric/null modular premises, not the matrix/fixed-point bookkeeping at one regulator scale.

\subsection{1.3 Edge-Center Completion (EC)}\label{13-edge-center-completion-ec}

\textbf{Theorem 1.1 (Derived EC decomposition).} Under the fixed-cutoff regulator realization of Section 1.2, and on the ordinary or central-defect branch where overlap rechartings along the cut reduce to a compact boundary group \(K_\Sigma\) or compact central extension \(\widehat K_\Sigma\), the collar Hilbert space is
\[
\mathcal{H}_{B_\delta}
=
\left(\tilde{\mathcal H}_{B_L}\otimes \tilde{\mathcal H}_{B_R}\right)^{\widehat K_\Sigma}
=
\bigoplus_{\alpha} \left(\mathcal{H}_{b_L^\alpha} \otimes \mathcal{H}_{b_R^\alpha}\right).
\]

The center of the collar algebra is generated by block projectors:
\[
Z(\mathcal{A}(B_\delta)) = \bigoplus_\alpha \mathbb{C} \cdot \mathbf{1}_\alpha.
\]

The proof is the Schur-lemma decomposition of the invariant subspace after decomposing the left and right boundary data into irreducible and contragredient \(\widehat K_\Sigma\)-modules. Thus EC is a theorem of overlap consistency plus the derived regulator package, not a model-specific extra premise.

\textbf{Corollary 1.2 (Exact Markov normal form as an extra state condition).} If, in addition, the reference collar state satisfies
\[
I_\omega(A_\delta : D_\delta \mid B_\delta) = 0,
\]
or one passes to an explicitly stated idealized recoverability limit that reduces to it, then
\[
\rho_{A_\delta B_\delta D_\delta}
=
\bigoplus_\alpha p_\alpha \left(\rho_{A_\delta b_L^\alpha} \otimes \rho_{b_R^\alpha D_\delta}\right).
\]

Approximate recoverability gives controlled deviations from this normal form; it is not implied by EC alone. The only fixed-cutoff obstruction to the ordinary-group form of Theorem 1.1 is a genuinely noncentral 2-group defect, which requires higher-gauge bookkeeping. The later open burden is the controlled refinement or scaling regime and the modular or null-limit premises, not the finite-dimensional collar decomposition itself.

\subsection{1.4 Fundamental Parameters}\label{14-fundamental-parameters}

OPH has exactly two fundamental screen parameters:

\begin{enumerate}
\def\labelenumi{\arabic{enumi}.}
\tightlist
\item
  \textbf{Pixel area:} \(a_{\text{cell}} \approx 1.63 \ell_P^2\)
\item
  \textbf{Screen capacity:} \(\log(\dim \mathcal{H}_{\text{tot}}) \sim 10^{122}\)
\end{enumerate}

From these:

\begin{itemize}
\tightlist
\item
  Newton\textquotesingle s constant: \(G = \frac{a_{\text{cell}}}{4\bar{\ell}(t)}\)
\item
  Cosmological constant: \(\Lambda = \frac{3\pi}{G \cdot \log(\dim \mathcal{H}_{\text{tot}})}\)
\end{itemize}

\begin{center}\rule{0.5\linewidth}{0.5pt}\end{center}

\section{2. Emergence of Gravity and Quantum Mechanics}\label{2-emergence-of-gravity-and-quantum-mechanics}

\subsection{2.1 Local Lorentz Structure}\label{21-local-lorentz-structure}

\textbf{Theorem 2.1 (Conformal Group Isomorphism):} The orientation-preserving conformal group of \(S^2\) is isomorphic to the connected Lorentz group: \[\text{Conf}^+(S^2) \cong \text{PSL}(2,\mathbb{C}) \cong \text{SO}^+(3,1)\]

This provides the kinematic structure of special relativity.

\subsection{2.2 Geometric Modular Flow}\label{22-geometric-modular-flow}

\textbf{Assumption G (Euclidean Regularity):} Modular flow near a smooth entangling cut has a regular Euclidean continuation, fixing the angular period to \(2\pi\).

\textbf{Theorem 2.2 (Bisognano-Wichmann on \(S^2\)):} In the refinement/scaling-limit EFT regime, under Assumption G and CMFP conditions, cap modular flow acts as the unique cap-preserving conformal dilation with period \(2\pi\): \[\sigma_t^\omega(A) = \Delta^{it} A \Delta^{-it}\]

The modular Hamiltonian takes the geometric form: \[K_C = 2\pi \int_C \beta(x) T_{00}(x) \, d^3x\] where \(\beta(x)\) is the local inverse temperature (Tolman factor).

\subsection{2.3 Null Modular Bridge (EFT Conditions N1-N3)}\label{23-null-modular-bridge-eft-conditions-n1-n3}

\textbf{Condition N1:} Consecutive null strips satisfy exact-or-controlled Markov modular additivity up to central boundary terms.

\textbf{Condition N2:} Null half-line algebras satisfy the half-sided modular inclusion property, either as an explicit scaling-limit premise or as the blow-up corollary of geometric modular flow near a smooth cut.

\textbf{Condition N3:} Bounded-interval modular Hamiltonians obey endpoint-Lipschitz control, and the renormalized null half-line family satisfies the weak differentiability / vanishing-at-infinity assumptions needed for the half-line endpoint-differentiation formula.

\subsection{2.4 Stress Tensor Reconstruction}\label{24-stress-tensor-reconstruction}

\textbf{Lemma 2.3 (Null Data Ambiguity):} If a symmetric tensor \(X_{ab}\) satisfies \(X_{ab}k^a k^b = 0\) for all null \(k\), then: \[X_{ab} = \phi \, g_{ab}\] for some scalar \(\phi\).

\emph{Proof:} Decompose \(X_{ab} = Y_{ab} + \phi g_{ab}\) with \(Y\) traceless. For null \(k = (1, \hat{n})\): \[X_{kk} = Y_{kk} + \phi g_{kk} = Y_{kk}\] since \(g_{kk} = 0\). Expanding: \[Y_{kk} = Y_{00} + 2\hat{n}^i Y_{0i} + \hat{n}^i \hat{n}^j Y_{ij}\]

If this vanishes for all \(\hat{n} \in S^2\), each angular moment (scalar, vector, traceless tensor) vanishes independently, so \(Y_{ab} = 0\). \(\square\)

\textbf{Corollary 2.4:} Null modular data determine \(T_{ab}\) only up to \(\phi g_{ab}\). Consequently, the Einstein equation is fixed only up to \(\Lambda g_{ab}\).

\textbf{Reconstruction Formulas:} From \(f(\hat{n}) := T_{kk}(\hat{n})\): \[T_{0i} = \frac{3}{2}\langle \hat{n}_i f \rangle_{S^2}\] \[T_{ij} - \frac{\delta_{ij}}{3}T^{(\text{spatial})}_{kk} = \frac{15}{2}\left\langle \left(\hat{n}_i \hat{n}_j - \frac{\delta_{ij}}{3}\right) f \right\rangle_{S^2}\]

\subsection{2.5 Entanglement Equilibrium and Einstein\textquotesingle s Equation}\label{25-entanglement-equilibrium-and-einsteins-equation}

\textbf{Theorem 2.5 (Conditional Jacobson-type Derivation):} In the same local Lorentzian scaling regime, under MaxEnt state selection with generalized entropy stationarity, the null-stress identification premise of the relativistic EFT branch, and admissible first variations about a maximally symmetric reference state,
\[\delta S_{\text{gen}} = 0\]

implies the rest-frame first-variation relation
\[
\delta\!\left(G_{00} + \Lambda g_{00}\right) = 8\pi G \,\delta\langle T_{00} \rangle.
\]

If this first-variation relation holds for all local directions and reference states in the scaling regime, overlap consistency upgrades it to the full tensor equation
\[
G_{ab} + \Lambda g_{ab} = 8\pi G \langle T_{ab} \rangle
\]
modulo the expected \(\Lambda g_{ab}\) ambiguity.

\textbf{Derivation sketch:}

\begin{enumerate}
\def\labelenumi{\arabic{enumi}.}
\item
  Raychaudhuri equation relates area change to \(R_{kk}\): \[\frac{d\theta}{d\lambda} = -\frac{1}{2}\theta^2 - \sigma^2 - R_{kk}\]
\item
  Recoverable Generalized Entropy implies: \[\frac{d}{d\lambda}\left(\frac{A}{4G} + S_{\text{bulk}}\right) \leq 0\]
\item
  Stationarity at MaxEnt gives: \[R_{kk} = 8\pi G \langle T_{kk} \rangle\]
\item
  By Lemma 2.3, this fixes Einstein\textquotesingle s equation up to \(\Lambda g_{ab}\).
\end{enumerate}

\subsection{2.6 Newton\textquotesingle s Constant from Edge Entropy}\label{26-newtons-constant-from-edge-entropy}

\textbf{Theorem 2.6:} The gravitational coupling is fixed by the edge entropy density: \[G = \frac{a_{\text{cell}}}{4\bar{\ell}(t)}\]

where \(\bar{\ell}(t) = a_{\text{cell}}/(4\ell_P^2)\) is the dimensionless entropy per cell.

With \(a_{\text{cell}} \approx 1.63 \ell_P^2\): \[\bar{\ell} = \frac{1.63}{4} \approx 0.4077\]

\begin{center}\rule{0.5\linewidth}{0.5pt}\end{center}

\section{3. The Measurement Problem}\label{3-the-measurement-problem}

\subsection{3.1 Records as Central Variables}\label{31-records-as-central-variables}

\textbf{Definition 3.1:} A \emph{record} in patch \(P\) is a family of orthogonal projectors \(\{\Pi_r\} \subset \mathcal{A}(P)\) such that:

\begin{enumerate}
\def\labelenumi{\arabic{enumi}.}
\tightlist
\item
  \textbf{Classicality:} \(\mathcal{R} := \text{vN}(\{\Pi_r\})\) is approximately abelian
\item
  \textbf{Stability:} \(\{\Pi_r\}\) are stable under effective dynamics
\item
  \textbf{Shareability:} \(\Pi_r\) lies in the center of overlap algebras
\end{enumerate}

\subsection{3.2 Markov Structure Theorem}\label{32-markov-structure-theorem}

\textbf{Theorem 3.1 (Quantum Markov Chain Structure):} If \(I(A:C|B) = 0\) for a tripartite state \(\rho_{ABC}\), then: \[\mathcal{H}_B \cong \bigoplus_j \left(\mathcal{H}_{b_L^{(j)}} \otimes \mathcal{H}_{b_R^{(j)}}\right)\]

and the state decomposes as: \[\rho_{ABC} = \bigoplus_j q_j \, \rho_{A b_L^{(j)}}^{(j)} \otimes \rho_{b_R^{(j)} C}^{(j)}\]

The label \(j\) is a \textbf{classical variable} living in the center of \(\mathcal{A}(B)\).

\subsection{3.3 Approximate Markov and Recovery}\label{33-approximate-markov-and-recovery}

\textbf{Theorem 3.2 (Fawzi-Renner Recovery):} If \(I(A:C|B) \leq \varepsilon\) (in bits), there exists a CPTP map \(\mathcal{R}: \mathcal{A}(B) \to \mathcal{A}(BC)\) such that: \[\left\|\rho_{ABC} - (\text{id}_A \otimes \mathcal{R})(\rho_{AB})\right\|_1 \leq 2\sqrt{\ln 2 \cdot \varepsilon}\]

\subsection{3.4 Definiteness of Outcomes}\label{34-definiteness-of-outcomes}

\textbf{Theorem 3.3 (Collapse from EC Structure):} Under EC, the post-measurement physical state is exactly a convex mixture over sector labels \(\alpha\): \[\rho = \bigoplus_\alpha p_\alpha \, \rho^{(\alpha)}\]

For any observer whose accessible observables lie in either side algebra:

\begin{enumerate}
\def\labelenumi{\arabic{enumi}.}
\tightlist
\item
  There are \textbf{no interference observables} between different \(\alpha\)-blocks
\item
  Each observer-instance has a definite \(\alpha\) in the classical sense
\end{enumerate}

\emph{Proof:} Cross-block operators \(|i\rangle\langle j|\) for \(i \in \alpha, j \in \alpha'\) (\(\alpha \neq \alpha'\)) are not gauge-invariant, hence not in the physical algebra. \(\square\)

\subsection{3.5 Born Rule as Unique Probability Measure}\label{35-born-rule-as-unique-probability-measure}

\textbf{Theorem 3.4 (Gleason):} For \(\dim \mathcal{H} \geq 3\), any measure \(\mu\) on projectors satisfying:

\begin{enumerate}
\def\labelenumi{\arabic{enumi}.}
\tightlist
\item
  Noncontextuality: \(\mu(P)\) depends only on \(P\)
\item
  Additivity: \(\mu(P \vee Q) = \mu(P) + \mu(Q)\) for \(P \perp Q\)
\item
  Normalization: \(\mu(\mathbf{1}) = 1\)
\end{enumerate}

must have the form: \[\mu(P) = \text{Tr}(\rho P)\] for some density operator \(\rho\).

\textbf{Corollary 3.5:} In OPH with the type-I regulator premise (R0), Born\textquotesingle s rule is the unique overlap-consistent probability assignment.

\subsection{3.6 Collapse as Conditioning}\label{36-collapse-as-conditioning}

\textbf{Proposition 3.6 (L\"{u}ders Update):} For record projectors \(\{\Pi_r\}\) and state \(\omega\): \[p_r = \omega(\Pi_r), \qquad \omega_r(X) = \frac{\omega(\Pi_r X \Pi_r)}{\omega(\Pi_r)}\]

Then for all \(X\) commuting with \(\mathcal{R}\): \[\omega(X) = \sum_r p_r \, \omega_r(X)\]

This is classical conditional probability on the record algebra.

\begin{center}\rule{0.5\linewidth}{0.5pt}\end{center}

\section{4. The Cosmological Principle}\label{4-the-cosmological-principle}

\subsection{4.1 MaxEnt Symmetry Inheritance}\label{41-maxent-symmetry-inheritance}

\textbf{Lemma 4.1:} Let \(G\) act on the screen by automorphisms \(\alpha_g\). If:

\begin{enumerate}
\def\labelenumi{\arabic{enumi}.}
\tightlist
\item
  The constraint set is \(G\)-invariant
\item
  Von Neumann entropy is \(G\)-invariant: \(S(\alpha_g(\rho)) = S(\rho)\)
\end{enumerate}

Then the MaxEnt optimizer \(\omega\) can be taken \(G\)-invariant. If unique, it is automatically \(G\)-invariant.

\emph{Proof:} If \(\omega\) is a maximizer, so is \(\omega \circ \alpha_{g^{-1}}\) by constraint invariance and entropy invariance. By uniqueness, \(\omega = \omega \circ \alpha_{g^{-1}}\). \(\square\)

\subsection{4.2 Isotropy from SO(3)-Invariant Constraints}\label{42-isotropy-from-so3-invariant-constraints}

\textbf{Assumption C:} Constraints are SO(3)-invariant on \(S^2\).

\textbf{Theorem 4.2:} Under Assumption C and MaxEnt uniqueness, the reference state \(\omega\) satisfies: \[\omega \circ \alpha_g = \omega \quad \forall g \in \text{SO}(3)\]

The stress tensor has perfect-fluid form: \[\langle T_{ab} \rangle = \rho \, u_a u_b + p(g_{ab} + u_a u_b)\]

\emph{Proof of perfect-fluid form:} Under SO(3), rank-2 tensors decompose into scalar + vector + traceless tensor irreps. Only scalars are invariant. Hence \(\langle T_{0i} \rangle = 0\) and \(\langle T_{ij} \rangle \propto \delta_{ij}\). \(\square\)

\subsection{4.3 Homogeneity from Isotropy Everywhere}\label{43-homogeneity-from-isotropy-everywhere}

\textbf{Theorem 4.3 (Schur-type):} Let \((\Sigma, h_{ij})\) be a 3-dimensional Riemannian manifold. If at every point \(p \in \Sigma\), the curvature tensor is SO(3)-invariant (pointwise isotropy), then: \[R_{ijkl}(p) = K(p)(h_{ik}h_{jl} - h_{il}h_{jk})\]

and by the second Bianchi identity, \(\nabla_m K = 0\), so \(K = \text{const}\).

\textbf{Corollary 4.4:} If OPH supplies isotropy for all observers, spatial geometry is a constant-curvature space form: \(S^3\), \(\mathbb{R}^3\), or \(H^3\).

\subsection{4.4 FLRW Emergence}\label{44-flrw-emergence}

\textbf{Theorem 4.5:} With:

\begin{enumerate}
\def\labelenumi{\arabic{enumi}.}
\tightlist
\item
  Semiclassical Einstein equation from recoverable generalized entropy / entanglement equilibrium
\item
  Perfect-fluid stress tensor from SO(3)-invariant MaxEnt
\item
  Positive \(\Lambda\) from finite screen capacity
\end{enumerate}

The emergent geometry is FLRW: \[ds^2 = -dt^2 + a(t)^2 \left(\frac{dr^2}{1-kr^2} + r^2 d\Omega^2\right)\]

with Friedmann equations: \[H^2 + \frac{k}{a^2} = \frac{8\pi G}{3}\rho + \frac{\Lambda}{3}\] \[\dot{H} - \frac{k}{a^2} = -4\pi G(\rho + p)\]

\begin{center}\rule{0.5\linewidth}{0.5pt}\end{center}

\section{5. Supersymmetry and Gauge Coupling Unification}\label{5-supersymmetry-and-gauge-coupling-unification}

\subsection{5.1 One-Loop Unification Algebra}\label{51-one-loop-unification-algebra}

At one loop, if couplings unify at \((M_U, \alpha_U)\): \[\alpha_i^{-1}(M_Z) = \alpha_U^{-1} + \frac{b_i}{2\pi}\ln\frac{M_U}{M_Z}\]

Define \(A_i := \alpha_i^{-1}(M_Z)\), \(L := \ln(M_U/M_Z)\). Eliminating \(\alpha_U\): \[L = \frac{2\pi}{b_1 - b_2}(A_1 - A_2)\]

The consistency relation for \(\alpha_s\) (a consistency check when all three couplings calibrate $P$; a genuine prediction in the two-input mode): \[A_3^{\text{pred}} = \frac{b_3 - b_2}{b_1 - b_2}A_1 + \frac{b_1 - b_3}{b_1 - b_2}A_2\]

\subsection{5.2 SM vs MSSM Coefficients}\label{52-sm-vs-mssm-coefficients}

{\def\LTcaptype{none} % do not increment counter
\begin{longtable}[]{@{}>{\RaggedRight\arraybackslash}p{0.320\textwidth}>{\RaggedRight\arraybackslash}p{0.320\textwidth}>{\RaggedRight\arraybackslash}p{0.320\textwidth}@{}}
\toprule\noalign{}
Model & \((b_1, b_2, b_3)\) & \(\alpha_s(M_Z)^{\text{pred}}\) \\
\midrule\noalign{}
\endhead
\bottomrule\noalign{}
\endlastfoot
SM & \((41/10, -19/6, -7)\) & \(\approx 0.071\) \\
MSSM & \((33/5, 1, -3)\) & \(\approx 0.116\) \\
Observed & , & \(0.1179 \pm 0.0010\) \\
\end{longtable}
}

The SM prediction is catastrophically wrong; MSSM-like coefficients work.

\subsection{5.3 Heat-Kernel Edge Sector Weights}\label{53-heat-kernel-edge-sector-weights}

\textbf{Theorem 5.1:} Under MaxEnt with bi-invariant constraints on a compact Lie group \(G\), the sector probabilities take heat-kernel form: \[p_R(t) \propto d_R \, e^{-t \, C_2(R)}\]

where \(d_R = \dim R\) and \(C_2(R)\) is the quadratic Casimir.

\textbf{Lemma 5.2 (Bi-invariant Operators):} If \(G=\prod_i G_i\) is a compact semisimple Lie group written as a product of compact simple factors, then any bi-invariant, second-order differential operator on \(G\) has the form: \[D = c_0 \mathbf{1} - \sum_i c_i \Delta_{G_i}\]

where \(\Delta_G\) is the Laplace-Beltrami operator.

\emph{Proof:} Bi-invariance \(\Leftrightarrow\) \(D \in Z(U(\mathfrak{g}))\) (center of universal enveloping algebra). For degree \(\leq 2\):

\begin{itemize}
\tightlist
\item
  Linear terms vanish (no invariant vectors in adjoint rep)
\item
  Quadratic terms are proportional to the Killing form on each simple factor, giving one independent coefficient per factor
\end{itemize}

The Casimir element acts as \(-\Delta_G\) in the regular representation. \(\square\)

\subsection{5.4 Peter-Weyl and Effective Multiplicity}\label{54-peter-weyl-and-effective-multiplicity}

By Peter-Weyl decomposition, the effective refinement-limit edge representation space is
\[L^2(G) \cong \bigoplus_R V_R \otimes V_R^*\]

\textbf{Key insight:} Entanglement traces over one factor give multiplicity \(d_R\) in \(p_R\), but loops see both factors, giving effective multiplicity: \[N_{\text{eff}}(R) = d_R \cdot p_R\]

\subsection{5.5 Beta Function Shifts}\label{55-beta-function-shifts}

The one-loop beta shift from edge sectors: \[\Delta b_a = \sum_R p_R \cdot d_R \cdot T_a(R)\]

where \(T_a(R)\) is the Dynkin index for gauge factor \(a\).

At the unification diffusion parameter \(t_U \approx 1.64\): \[\Delta b \approx (2.49, 4.38, 3.97) \quad \text{vs MSSM} \quad (2.50, 4.17, 4.00)\]

With fermionic grading (restricting to half-integer SU(2) sectors): \[\Delta b \approx (2.50, 4.17, 3.97)\]

matching MSSM within 1\%.

\begin{center}\rule{0.5\linewidth}{0.5pt}\end{center}

\section{6. The Horizon Problem}\label{6-the-horizon-problem}

\subsection{6.1 Anisotropy Bound from Markov Control}\label{61-anisotropy-bound-from-markov-control}

\textbf{Theorem 6.1:} For a tripartition \(A\)-\(B\)-\(D\) with collar width \(\delta\), the Markov error satisfies: \[I(A:D|B) \lesssim c \cdot |\partial C|_{\text{UV}} \cdot e^{-\delta/\xi}\]

where \(\xi\) is the correlation length.

\textbf{Corollary 6.2:} For bounded observable \(O\) with \(\|O\| \sim 1\): \[|\Delta\langle O \rangle| \leq 2\sqrt{\ln 2 \cdot \varepsilon}\]

\subsection{6.2 CMB Anisotropy Requirement}\label{62-cmb-anisotropy-requirement}

For \(\delta T/T \lesssim 10^{-5}\): \[2\sqrt{\ln 2 \cdot \varepsilon} \lesssim 10^{-5}\] \[\varepsilon \lesssim 3.61 \times 10^{-11} \text{ bits}\]

Required collar width: \[\frac{\delta}{\xi} \gtrsim \ln\left(\frac{c \cdot |\partial C|_{\text{UV}}}{\varepsilon_{\text{max}}}\right) \approx 24\]

With \(\xi \approx \sqrt{a_{\text{cell}}} \cdot \ell_P \approx 1.28 \ell_P\): \[\delta_{\text{CMB}} \approx 31 \ell_P\]

\subsection{6.3 Homogeneity as MaxEnt Default}\label{63-homogeneity-as-maxent-default}

\textbf{Theorem 6.3:} If MaxEnt constraints are uniform across UV cells with no marked points, the MaxEnt state is invariant under triangulation automorphisms. In the continuum limit, this gives SO(3) invariance.

\emph{Proof:} Standard MaxEnt uniqueness + symmetry inheritance (Lemma 4.1). \(\square\)

\begin{center}\rule{0.5\linewidth}{0.5pt}\end{center}

\section{7. Black Hole Information Paradox}\label{7-black-hole-information-paradox}

\subsection{7.1 The Information Trilemma}\label{71-the-information-trilemma}

The AMPS-style paradox assumes:

\begin{enumerate}
\def\labelenumi{\arabic{enumi}.}
\tightlist
\item
  Outgoing mode \(B\) entangled with interior partner \(A\) (smooth horizon)
\item
  Late radiation \(B\) entangled with early radiation \(R\) (unitarity)
\item
  Monogamy: \(B\) cannot be maximally entangled with both
\end{enumerate}

\textbf{The false assumption:} \(\mathcal{H} \stackrel{?}{=} \mathcal{H}_{\text{inside}} \otimes \mathcal{H}_{\text{outside}} \otimes \mathcal{H}_{\text{radiation}}\)

\subsection{7.2 Edge-Center Resolution}\label{72-edge-center-resolution}

\textbf{Theorem 7.1:} Under EC, the collar decomposition gives: \[\rho_{A_\delta B_\delta D_\delta} = \bigoplus_\alpha p_\alpha \left(\rho_{A_\delta b_L^\alpha} \otimes \rho_{b_R^\alpha D_\delta}\right)\]

The "glue" between inside and outside is the edge sector label \(\alpha\) in the center. Conditional on \(\alpha\), inside and outside factorize.

\subsection{7.3 Recoverability of Interior}\label{73-recoverability-of-interior}

By the recoverability clause of the canonical entropy axiom (small CMI) + Theorem 3.2: \[\left\|\rho_{ABC} - (\text{id}_A \otimes \mathcal{R})(\rho_{AB})\right\|_1 \leq 2\sqrt{\ln 2 \cdot \varepsilon}\]

\textbf{Interpretation:} The "interior partner" is recoverable from outside + radiation. It\textquotesingle s encoded, not independent.

\subsection{7.4 Hawking Temperature from Modular Theory}\label{74-hawking-temperature-from-modular-theory}

\textbf{Theorem 7.2:} The outside algebra is in a KMS state at inverse temperature \(\beta = 2\pi/\kappa\) with respect to the horizon Killing generator.

\emph{Derivation:} Near the horizon: \[ds^2 \approx -\kappa^2 \rho^2 dt^2 + d\rho^2 + r_H^2 d\Omega^2\]

Euclidean regularity at \(\rho = 0\) requires \(\kappa t_E \sim \kappa t_E + 2\pi\), hence: \[T_H = \frac{\kappa}{2\pi} = \frac{\hbar c^3}{8\pi G k_B M}\]

\subsection{7.5 Discrete Area Spectrum}\label{75-discrete-area-spectrum}

From the central area operator: \[L_C = \sum_\alpha (\log d_\alpha) P_\alpha\]

Area eigenvalues: \(A_\alpha = 4G \log d_\alpha = 4\ell_P^2 \ln d_\alpha\)

For Schwarzschild, a sector transition \(d \to d'\) gives: \[\Delta E = k_B T_H \ln(d'/d)\]

\begin{center}\rule{0.5\linewidth}{0.5pt}\end{center}

\section{8. The Cosmological Constant Problem}\label{8-the-cosmological-constant-problem}

\subsection{8.1 Vacuum Energy Blindness}\label{81-vacuum-energy-blindness}

\textbf{Lemma 8.1:} For any null vector \(k\): \[T^{\text{vac}}_{kk} = T^{\text{vac}}_{ab} k^a k^b = -\rho_{\text{vac}} g_{ab} k^a k^b = -\rho_{\text{vac}}(k \cdot k) = 0\]

Vacuum energy contributions lie in the kernel of the null map \(T_{ab} \mapsto T_{kk}\).

\subsection{8.2 Structural Separation}\label{82-structural-separation}

\textbf{Proposition 8.2:} In OPH:

\begin{itemize}
\tightlist
\item
  Local modular/null data fixes \(T_{ab}\) up to \(\phi g_{ab}\)
\item
  \(\Lambda\) is fixed by global screen capacity, not local physics
\end{itemize}

The "huge zero-point energy" simply isn\textquotesingle t the thing determining curvature.

\subsection{8.3 $\Lambda$ from Capacity}\label{83-ux3bb-from-capacity}

\[\Lambda = \frac{3\pi}{G \cdot \log(\dim \mathcal{H}_{\text{tot}})}\]

Equivalently, using de Sitter entropy: \[S_{dS} = \frac{A_{dS}}{4G} = \frac{3\pi}{G\Lambda} = \log(\dim \mathcal{H}_{\text{tot}})\]

For observed \(\Lambda \approx 1.09 \times 10^{-52}\) m\(^{-2}\): \[\log(\dim \mathcal{H}_{\text{tot}}) \sim 10^{122}\]

\subsection{8.4 No-Go for Local $\Lambda$ Prediction}\label{84-no-go-for-local-ux3bb-prediction}

\textbf{Theorem 8.3:} Within OPH\textquotesingle s null-modular reconstruction:

\begin{enumerate}
\def\labelenumi{\arabic{enumi}.}
\tightlist
\item
  All \(T_{kk}\) data unchanged under \(T_{ab} \to T_{ab} + \phi g_{ab}\)
\item
  Therefore \(\Lambda\) cannot be fixed by local overlap consistency
\item
  \(\Lambda\) requires a global input (capacity)
\end{enumerate}

\begin{center}\rule{0.5\linewidth}{0.5pt}\end{center}

\section{9. Dark Matter as Modular Anomaly}\label{9-dark-matter-as-modular-anomaly}

\subsection{9.1 The Modular Additivity Defect}\label{91-the-modular-additivity-defect}

Define the collar modular additivity defect: \[\Delta K_\delta := K_{ABD} - K_{AB} - K_{BD} + K_B\]

\textbf{Theorem 9.1:} \(\langle \Delta K_\delta \rangle_\omega = -I(A:D|B)_\omega\)

This is the conditional mutual information, measuring Markov imperfection.

\subsection{9.2 Anomalous Stress-Energy}\label{92-anomalous-stress-energy}

The modified Einstein equation: \[G_{ab} + \Lambda g_{ab} = 8\pi G \left(\langle T_{ab} \rangle + \langle T_{ab}^{\text{anom}} \rangle\right)\]

where in the diamond rest frame: \[\langle T_{00}^{\text{anom}} \rangle := \frac{15}{8\pi^2} \frac{\delta\langle K_C^{(\text{anom})} \rangle}{\ell^4}\]

The coefficient \(15/(8\pi^2)\) comes from inverting the \(d=4\) geometric integral coefficient \(\Omega_2/(4^2-1) = 4\pi/15\).

\subsection{9.3 Properties of the Anomaly Sector}\label{93-properties-of-the-anomaly-sector}

\begin{enumerate}
\def\labelenumi{\arabic{enumi}.}
\tightlist
\item
  \textbf{Gravitates:} Appears on RHS of Einstein equation
\item
  \textbf{Dark by construction:} Encoded in central/recoverability data, not SM fields
\item
  \textbf{Covariantly conserved:} \(\nabla^a T_{ab}^{\text{anom}} = 0\) (from Bianchi identity)
\item
  \textbf{Effectively classical:} Central structure implies classical labels
\end{enumerate}

\subsection{9.4 The MOND Acceleration Scale}\label{94-the-mond-acceleration-scale}

\textbf{Corollary 9.2 (program-level):} The unique IR acceleration scale available to the modular-anomaly branch from \(\Lambda\) is \[a_0^{(\text{OPH})} := \frac{15}{8\pi^2} c^2 \sqrt{\frac{\Lambda}{3}} = \frac{15}{8\pi^2} \frac{c^2}{r_{dS}}\]

\textbf{Numerical evaluation:} \[a_0^{(\text{OPH})} \approx 1.03 \times 10^{-10} \text{ m/s}^2\]

Compare to observed: \(a_0^{(\text{obs})} \approx 1.2 \times 10^{-10}\) m/s\(^2\) (within 15\%).

\subsection{9.5 Polarization Response (MOND-like Scaling)}\label{95-polarization-response-mond-like-scaling}

\textbf{Uniqueness argument:} In the deep-IR regime, the only scales are \(a_0\) and \(g_b\) (Newtonian baryonic acceleration). Dimensional analysis + scale invariance forces: \[g_{\text{anom}} \propto \sqrt{a_0 \cdot g_b}\]

The cubic polarization functional: \[S_{\text{pol}}[\varphi] = \int d^3x \left[-\frac{1}{12\pi G a_0}|\nabla\varphi|^3 - \rho_b \varphi\right]\]

yields the MOND/RAR form: \[g_{\text{obs}} \approx g_b + \sqrt{a_0 g_b}\]

\begin{center}\rule{0.5\linewidth}{0.5pt}\end{center}

\section{10. Proton Stability and Spin}\label{10-proton-stability-and-spin}

\subsection{10.1 Sector Factorization and Product Group}\label{101-sector-factorization-and-product-group}

\textbf{Theorem 10.1 (Factorization Equivalence):} \[\text{Factorizing edge weights} \Leftrightarrow \text{Additive boundary Laplacian} \Leftrightarrow \text{Product gauge group}\]

If \(H_\partial = H_\partial^{(1)} + H_\partial^{(2)} + H_\partial^{(3)}\) with \([H_\partial^{(i)}, H_\partial^{(j)}] = 0\): \[p(R_1, R_2, R_3) \propto \prod_{i=1}^3 d_{R_i} e^{-t_i C_2(R_i)}\]

Applied to the refinement-stable edge-sector colimit, Tannaka-Krein reconstruction first yields some compact group \(G\). Under the extended MAR admissibility package, and once the admissible class includes one connected abelian charge factor, the connected Lie realization selected by MAR is \(SU(3) \times SU(2) \times U(1)\) (up to finite quotient).

Under the extended theory \(T_{\mathrm{ext}}\), the explicit MAR selector picks this product structure on the admissible class: the minimal coupled carrier \(\mathbb{C}^3 \otimes \mathbb{C}^2\) enforces commuting color and weak actions, which implies the additive boundary Laplacian. See Part II of this manuscript for the complete proof.

\subsection{10.2 No Leptoquarks, No Gauge Proton Decay}\label{102-no-leptoquarks-no-gauge-proton-decay}

\textbf{Corollary 10.2:} With product gauge group, the adjoint representation of the full connected gauge group is \[(8,1,0) \oplus (1,3,0) \oplus (1,1,0)\]

equivalently, the adjoint of the derived gauge group is
\[
(8,1,0) \oplus (1,3,0).
\]

No gauge generators in mixed representations \((3,2,\pm 5/6)\) (the SU(5) \(X,Y\) bosons). Hence: \[\tau_p^{(\text{gauge})} = \infty\]

\subsection{10.3 Proton Spin Fraction}\label{103-proton-spin-fraction}

\textbf{Casimir Equilibration Ansatz:} Spin sharing between quarks and gluons equilibrates according to coupling strength: \[\Delta\Sigma \approx \frac{C_F}{C_F + C_A}\]

For SU(3): \(C_F = 4/3\), \(C_A = 3\): \[\Delta\Sigma = \frac{4/3}{4/3 + 3} = \frac{4}{13} \approx 0.308\]

Compare to lattice: \(\Delta\Sigma \approx 0.286\) (within 8\%).

\begin{center}\rule{0.5\linewidth}{0.5pt}\end{center}

\section{11. Baryogenesis}\label{11-baryogenesis}

\subsection{11.1 The Z$_6$ Defect Suppression}\label{111-the-zux2086-defect-suppression}

The SM global gauge structure is \((G_{\text{SM}})/\mathbb{Z}_6\). The entropy deficit from the \(\mathbb{Z}_6\) quotient: \[\Delta S = \ln 6\]

Defect suppression factor: \[\varepsilon = e^{-\Delta S} = \frac{1}{6}\]

\subsection{11.2 The Baryon-Violating Channel}\label{112-the-baryon-violating-channel}

The electroweak topological transition (\(\Delta N_{CS} = \pm 1\)) produces chiral zero modes. By index theorem:

\begin{itemize}
\tightlist
\item
  One zero mode per left-handed SU(2) doublet per unit topological charge
\item
  Per generation: \(N_c + 1 = 4\) doublets
\item
  Total: \(n_{\text{doublets}} = (N_c + 1)N_g = 12\)
\end{itemize}

The \textquotesingle t Hooft vertex is a 12-fermion interaction.

\subsection{11.3 Defect-Order Identification}\label{113-defect-order-identification}

\textbf{Proposition 11.1:} The baryon-violating channel corresponds to defect order \(n = 12\): \[\varepsilon^{12} = 6^{-12} \approx 4.59 \times 10^{-10}\]

Compare to observed \(\eta_B \approx 6.1 \times 10^{-10}\) (same order of magnitude).

\subsection{11.4 CP Bias from Overlap Holonomy}\label{114-cp-bias-from-overlap-holonomy}

Under CP: \(U_{ij} \mapsto U_{ij}^*\), \(\Omega_{ijk} \mapsto \Omega_{ijk}^\dagger\)

For central cocycle \(z_{ijk} \in U(1)\): CP sends \(z \to z^{-1}\).

\textbf{Real cocycle condition:} \([z] = [z]^{-1}\) (cohomologous to inverse)

Generic phases violate this \(\Rightarrow\) CP violation is generic in OPH.

\begin{center}\rule{0.5\linewidth}{0.5pt}\end{center}

\section{12. Generation Structure and Yukawa Hierarchy}\label{12-generation-structure-and-yukawa-hierarchy}

\subsection{12.1 Why Three Generations}\label{121-why-three-generations}

\textbf{Theorem 12.1 (N\_g = 3 Selection):}

\begin{enumerate}
\def\labelenumi{\arabic{enumi}.}
\tightlist
\item
  \textbf{CP violation requires N\_g ≥ 3:} Physical CKM phase exists only for \(N_g \geq 3\)
\item
  \textbf{UV stability bounds N\_g ≤ 5:} SU(2) asymptotic freedom
\item
  \textbf{Minimality selects N\_g = 3:} Extra generations introduce unconstrained relevant deformations
\end{enumerate}

\subsection{12.2 Yukawa as Defect-Mediated Overlap}\label{122-yukawa-as-defect-mediated-overlap}

\textbf{Theorem 12.2:} If a Yukawa coupling requires \(n\) units of \(\mathbb{Z}_6\) defect insertion: \[y \propto \langle D^n \rangle = 6^{-n}\]

The suppression factor is fixed by topology, not tuned.

\subsection{12.3 The Yukawa Spectrum}\label{123-the-yukawa-spectrum}

{\def\LTcaptype{none} % do not increment counter
\begin{longtable}[]{@{}>{\RaggedRight\arraybackslash}p{0.240\textwidth}>{\RaggedRight\arraybackslash}p{0.240\textwidth}>{\RaggedRight\arraybackslash}p{0.240\textwidth}>{\RaggedRight\arraybackslash}p{0.240\textwidth}@{}}
\toprule\noalign{}
Fermion & Yukawa & \(n = -\ln y / \ln 6\) & Nearest integer \\
\midrule\noalign{}
\endhead
\bottomrule\noalign{}
\endlastfoot
\(t\) & 0.992 & 0.00 & 0 \\
\(b\) & 0.024 & 2.08 & 2 \\
\(c\) & 0.0073 & 2.75 & 3 \\
\(s\) & 5.3×10⁻⁴ & 4.21 & 4 \\
\(d\) & 2.7×10⁻⁵ & 5.87 & 6 \\
\(u\) & 1.3×10⁻⁵ & 6.29 & 6 \\
\(\tau\) & 0.010 & 2.56 & 3 \\
\(\mu\) & 6.1×10⁻⁴ & 4.13 & 4 \\
\(e\) & 2.9×10⁻⁶ & 7.11 & 7 \\
\end{longtable}
}

The base-6 logarithm lands near integers across the spectrum.

\subsection{12.4 Deterministic Completion (RK Hamiltonian)}\label{124-deterministic-completion-rk-hamiltonian}

Define configuration weight: \[w(\mathcal{C}) = \exp\left(-t \sum_{\text{plaquettes}} C_2(R_p)\right) \times 6^{-n_D(\mathcal{C})}\]

The Rokhsar-Kivelson ground state: \[|\Psi\rangle = \frac{1}{\sqrt{Z}} \sum_{\mathcal{C}} \sqrt{w(\mathcal{C})} |\mathcal{C}\rangle\]

is the unique ground state of a frustration-free parent Hamiltonian, replacing MaxEnt ensemble with a deterministic pure state.

\begin{center}\rule{0.5\linewidth}{0.5pt}\end{center}

\section{13. The Koide Formula}\label{13-the-koide-formula}

\subsection{13.1 Numerical Verification}\label{131-numerical-verification}

Using PDG masses:

\begin{itemize}
\tightlist
\item
  \(m_e = 0.51099895\) MeV
\item
  \(m_\mu = 105.6583755\) MeV
\item
  \(m_\tau = 1776.93\) MeV
\end{itemize}

\[Q = \frac{m_e + m_\mu + m_\tau}{(\sqrt{m_e} + \sqrt{m_\mu} + \sqrt{m_\tau})^2} = 0.6666644645\]

Compare to \(2/3 = 0.6666666667\). Difference: \(2.2 \times 10^{-6}\) (within 1σ).

\subsection{13.2 Z$_3$ Holonomy Parameterization}\label{132-zux2083-holonomy-parameterization}

\textbf{Theorem 13.1:} The minimal Hermitian circulant on generation space: \[\Phi = a \cdot I + b \cdot P + b^* \cdot P^2\]

with \(P\) the cyclic shift, has eigenvalues: \[\lambda_k = a + 2|b|\cos\left(\delta + \frac{2\pi k}{3}\right)\]

where \(\delta = \arg(b)\).

\subsection{13.3 Why Q = 2/3}\label{133-why-q--23}

\textbf{Theorem 13.2:} With \(r_k = \lambda_k\) (root masses): \[Q = \frac{\sum_k r_k^2}{(\sum_k r_k)^2} = \frac{1 + 2(|b|/a)^2}{3}\]

Setting \(Q = 2/3\): \[|b|/a = \frac{1}{\sqrt{2}}\]

This is the "balanced" configuration where singlet and charged modes have equal norm.

\subsection{13.4 The Koide Phase from OPH}\label{134-the-koide-phase-from-oph}

\textbf{Proposition 13.3:} The holonomy phase: \[\delta_{\text{OPH}} = \frac{\beta_{\text{EW}} \cdot Y_Q}{N_g} = \frac{(N_c + 1)}{2N_c \cdot N_g}\]

With \(N_c = 3\), \(N_g = 3\): \[\delta_{\text{OPH}} = \frac{4}{18} = \frac{2}{9} = 0.2222...\]

Experimental extraction: \(\delta_{\text{exp}} = 0.2222248 \pm 0.0000063\)

Agreement within 0.4σ.

\begin{center}\rule{0.5\linewidth}{0.5pt}\end{center}

\section{14. Cosmological Parameters}\label{14-cosmological-parameters}

\subsection{14.1 The de Sitter Scale}\label{141-the-de-sitter-scale}

From \(\Lambda \approx 1.09 \times 10^{-52}\) m\(^{-2}\): \[r_{dS} = \sqrt{\frac{3}{\Lambda}} \approx 1.66 \times 10^{26} \text{ m}\]

The de Sitter time: \[t_\Lambda = \frac{r_{dS}}{c} \approx 5.53 \times 10^{17} \text{ s} \approx 17.5 \text{ Gyr}\]

\subsection{14.2 Emergent Age of the Universe}\label{142-emergent-age-of-the-universe}

For flat ΛCDM: \[t_0 = \frac{2}{3H_0\sqrt{\Omega_\Lambda}} \sinh^{-1}\left(\sqrt{\frac{\Omega_\Lambda}{\Omega_m}}\right)\]

With \(H_0 \approx 67.4\) km/s/Mpc, \(\Omega_\Lambda \approx 0.685\): \[t_0 \approx 13.8 \text{ Gyr}\]

\subsection{14.3 Screen Capacity}\label{143-screen-capacity}

\[\dim \mathcal{H}_{\text{tot}} = \exp(S_{dS}) = \exp\left(\frac{3\pi}{G\Lambda}\right)\]

\[\log(\dim \mathcal{H}_{\text{tot}}) \approx 2.85 \times 10^{122}\]

\subsection{14.4 Pixel Parameters}\label{144-pixel-parameters}

\begin{itemize}
\tightlist
\item
  Cell area: \(a_{\text{cell}} \approx 1.63 \ell_P^2\)
\item
  UV length: \(\ell_{\text{UV}} = \sqrt{a_{\text{cell}}} \approx 1.28 \ell_P\)
\item
  Entropy per cell: \(\bar{\ell} = a_{\text{cell}}/(4\ell_P^2) \approx 0.408\)
\end{itemize}

\begin{center}\rule{0.5\linewidth}{0.5pt}\end{center}

\section{Appendix A: Summary of Derived Results}\label{appendix-a-summary-of-derived-results}

\subsection{A.1 Theorem-Level Results and Dependency Status}\label{a1-theorem-level-results}

This appendix records theorem-level status only. The recovered core of the current OPH program is D1--D5 together with D7--D9. D6 is included below only as the adjacent Phase-II boundary marker; D10--D12 are excluded from the theorem-level list except where their non-core status must be made explicit. The table matches the compact note's authoritative ledger and keeps imported mathematics, extra premises, and claim tier visible.

\begingroup
\scriptsize
{\def\LTcaptype{none}
\setlength{\tabcolsep}{2pt}
\begin{longtable}[]{@{}>{\RaggedRight\arraybackslash}p{0.07\textwidth}>{\RaggedRight\arraybackslash}p{0.21\textwidth}>{\RaggedRight\arraybackslash}p{0.15\textwidth}>{\RaggedRight\arraybackslash}p{0.22\textwidth}>{\RaggedRight\arraybackslash}p{0.17\textwidth}@{}}
\toprule\noalign{}
Node & Output & Standard mathematics used & Extra premises / inputs beyond the five axioms & Status \\
\midrule\noalign{}
\endhead
\bottomrule\noalign{}
\endlastfoot
\textbf{D1} & schedule-independent normal form for overlap repair & Newman's lemma & finite patch net, termination, and local confluence & Phase I structural theorem \\
\textbf{D2} & EC decomposition, Markov collar, and generalized-entropy split & HJPW; Petz; Fawzi--Renner & R0/R1 or explicit recoverability control; coarse-grained \(A/(4G)\) dictionary & Phase I structural lemma/proposition layer \\
\textbf{D3} & geometric modular flow and Lorentz branch & Bisognano--Wichmann modular geometry; conformal-group classification & scaling-limit geometric modular phase and Euclidean regularity & Phase I conditional scaling-limit theorem \\
\textbf{D4} & null modular bridge to \(T_{kk}\) & Borchers--Wiesbrock; distributional differentiation on half-lines & null-strip control, half-sided inclusion, and the relativistic null-stress identification premise & Phase I conditional bridge theorem \\
\textbf{D5} & Jacobson-type Einstein branch & Jacobson small-ball mathematics; null-to-tensor reconstruction modulo a metric term & fixed-cap stationarity and the locally Lorentzian \(d=4\) scaling regime & Phase I conditional scaling-limit theorem \\
\textbf{D6} & capacity/\(\Lambda\) corollary and the separate capacity-based neutrino estimate & de Sitter entropy relation and dimensional analysis & external input \(N_{\mathrm{scr}}\) and the screen-capacity identification & Phase II input-dependent corollary / estimate \\
\textbf{D7} & compact gauge reconstruction in the bosonic branch & Doplicher--Roberts / Tannaka reconstruction & rigid symmetric bosonic refinement-stable sector colimit, transportability, and fiber-functor premises & Phase I conditional structural theorem \\
\textbf{D8} & product structure up to finite quotient on the realized branch & compact Lie representation classification; Schur's lemma & MAR, connected admissible class, one connected abelian factor, and \([z]=0\) & Phase I realized-branch theorem \\
\textbf{D9} & exact hypercharges, \(N_c=3\), \(N_g=3\), and the \(\mathbb Z_6\) quotient & anomaly algebra; Witten anomaly; CKM CP counting & realized one-generation/one-Higgs package and MAR admissibility premises & Phase I realized-branch theorem/corollary chain \\
\end{longtable}
}
\endgroup

The recovered core in this appendix is therefore D1--D5 together with D7--D9. D6 appears only to keep the Phase-I / Phase-II boundary explicit. Nothing in this appendix should be read as promoting D10, D11, or D12 to theorem level.
\subsection{A.2 Secondary Quantitative and Continuation Benchmarks}\label{a2-quantitative-predictions}

This appendix does not extend the theorem-level list. It records every non-core branch that is still useful to track:

\begin{enumerate}
\item \textbf{Phase II input-dependent corollary (D6):} the capacity relation for \(\Lambda\) and the separate order-of-magnitude neutrino estimate tied to the same capacity input.
\item \textbf{Phase II current implementation: calibration sector (D10):} gauge/electroweak closure of the present supplement-backed implementation.
\item \textbf{Phase II current implementation: secondary quantitative branch (D11):} the Higgs/top critical-surface branch once the D10 trajectory and the extra critical-surface condition are imposed.
\item \textbf{Phase III deferred continuations (D12 and beyond):} charged-lepton/Koide, texture, neutrino mass/mixing constructions beyond the D6 estimate, dark-sector, baryogenesis, black-hole spectroscopy, string/worldsheet, proton-spin, and other downstream fits.
\end{enumerate}

Failure of item 4 retracts the corresponding continuation only. Failure of item 3 challenges the D11 branch; failure of item 2 challenges the D10 implementation; failure of item 1 challenges the screen-capacity identification only. None of those failures by itself erases the recovered core D1--D5 plus D7--D9.

\subsection{A.3 Future Directions}\label{a3-future-directions}

\begin{enumerate}
\item construct or prove the scaling-limit regime required by the D3--D5 relativity chain, including the null-stress bridge and fixed-cap stationarity premises;
\item extend the compact-gauge reconstruction from the bosonic internal-gauge branch to the fermionic/super-Tannakian chiral branch while preserving the D7--D9 Standard Model structural chain;
\item derive or independently fix the external continuous inputs $P$ and $N_{\mathrm{scr}}$, and complete a fully auditable D10--D11 numerical pipeline with explicit running, threshold, and matching conventions;
\item separate the D6 capacity relation from flavor-style add-ons by deciding exactly which neutrino or cosmological continuations survive without additional ans\"atze;
\item only after items 1--4 are secured, reopen the Phase-III continuations: charged-lepton/Koide structure, texture exponents, dark-sector response laws, baryogenesis, black-hole spectroscopy, and string/worldsheet reorganization, including a controlled large-\(N_{\mathrm{edge}}\) regime distinct from \(N_c=3\).
\end{enumerate}

\begin{center}\rule{0.5\linewidth}{0.5pt}\end{center}

\section{Appendix B: Mixed-Status Numerical Benchmarks}\label{appendix-b-key-numerical-values}

{\def\LTcaptype{none} % do not increment counter
\begin{longtable}[]{@{}>{\RaggedRight\arraybackslash}p{0.240\textwidth}>{\RaggedRight\arraybackslash}p{0.240\textwidth}>{\RaggedRight\arraybackslash}p{0.240\textwidth}>{\RaggedRight\arraybackslash}p{0.240\textwidth}@{}}
\toprule\noalign{}
Quantity / branch & OPH Value & Observed & Status \\
\midrule\noalign{}
\endhead
\bottomrule\noalign{}
\endlastfoot
\(a_0\) (dark-sector continuation) & \(1.03 \times 10^{-10}\) m/s² & \(1.2 \times 10^{-10}\) m/s² & deferred continuation benchmark \\
\(\Delta\Sigma\) (spin continuation) & 0.308 & 0.29 ± 0.03 & deferred continuation benchmark \\
\(\eta_B\) (baryogenesis estimate) & \(4.6 \times 10^{-10}\) & \(6.1 \times 10^{-10}\) & deferred continuation benchmark \\
Koide \(Q\) (flavor continuation) & 2/3 & 0.666664 & deferred continuation benchmark \\
Koide \(\delta\) (flavor continuation) & 2/9 & 0.222225 & deferred continuation benchmark \\
\(\Delta b_3/\Delta b_2\) (D10 calibration) & 0.91 & 0.96 (MSSM) & supplement-backed calibration \\
\end{longtable}
}

\begin{center}\rule{0.5\linewidth}{0.5pt}\end{center}

\section{References}\label{references}

The derivations in this supplement are based on the OPH framework as specified in Part I of this manuscript. Standard results from:

\begin{itemize}
\tightlist
\item
  Tomita-Takesaki modular theory
\item
  Quantum Markov chain structure theorems (Hayden-Jozsa-Petz-Winter)
\item
  Fawzi-Renner recovery bounds
\item
  Gleason\textquotesingle s theorem
\item
  Jacobson\textquotesingle s entanglement equilibrium
\item
  Peter-Weyl decomposition
\item
  Heat kernel asymptotics on compact Lie groups
\item
  Tannaka-Krein reconstruction
\item
  \textquotesingle t Hooft anomaly matching and instanton calculus
\end{itemize}
